% =====================================================================
%  2bat-ud6-5.tex  —  SESSIÓ 5: Què vol dir ±3 per cent?
%  (sense preàmbul: s'inclou des de main.tex)
% =====================================================================
\horatitol{Sessió 5 — Què vol dir $\pm 3$ per cent?}%
{Activitat 2: descobrir, a partir de la fitxa tècnica d'una enquesta real, que el resultat no és un valor exacte sinó un interval de confiança.}

\subsection{Idea clau}

A les \textbf{pissarres}, en grups, analitzem la fitxa tècnica d'una enquesta real que
dona un resultat amb el seu \textbf{marge d'error} (per exemple, «el $42\%$ opina X, amb un
marge de $\pm 3\%$»). La pregunta que debatreu: què afirma \emph{realment} l'enquesta? Pot
dir que el percentatge exacte és $42\%$? Què aporta el «$\pm 3\%$»?

\begin{idea}
\textbf{L'interval de confiança.} Una enquesta no dona un valor exacte, sinó un
\textbf{interval} dins del qual és molt probable que estigui el valor real de tota la
població:
\[
\text{interval} = \text{resultat} \pm \text{marge d'error}.
\]
\begin{itemize}[leftmargin=1.4em, itemsep=2pt]
  \item Amb «$42\%\pm 3\%$», l'interval de confiança va del \textbf{39\%} al \textbf{45\%}.
  \item Vol dir: és molt probable que el valor real de la població estigui \emph{entre el
        39\% i el 45\%}, no que sigui exactament el $42\%$.
  \item El marge d'error és l'expressió \textbf{honesta} de la incertesa d'una estimació.
\end{itemize}
\textbf{Errors d'interpretació típics:} creure que el $42\%$ és exacte, o pensar que el
marge vol dir que «alguns han respost malament» (no té res a veure).
\end{idea}

\subsection{Exemples}

\begin{exemple}[del resultat a l'interval]
Una enquesta diu: «el $42\%$ de la població està a favor d'ampliar un servei, amb un marge
d'error de $\pm 3\%$». L'interval de confiança és
\[
42\%-3\%=39\% \quad\text{fins a}\quad 42\%+3\%=45\%.
\]
La lectura correcta és: «és molt probable que, a tota la població, el suport estigui
\textbf{entre el 39\% i el 45\%}». No es pot afirmar que sigui exactament el $42\%$.
\end{exemple}

\begin{exemple}[què NO vol dir el marge]
El «$\pm 3\%$» \textbf{no} vol dir que un $3\%$ de la gent s'hagi equivocat, ni que un $3\%$
menteixi. Vol dir que, com que s'ha preguntat a una mostra i no a tothom, el valor real pot
diferir del $42\%$ mesurat fins a un $3\%$ amunt o avall. És incertesa del mostreig, no
error de les persones.
\end{exemple}

\subsection{Exercicis resolts}

\begin{exresolt}[construir i interpretar l'interval]
Una enquesta dona «el $55\%$ de les famílies està satisfet, amb un marge de $\pm 4\%$».
Construeix l'interval i escriu què afirma correctament l'enquesta.
\solucio
Interval: $55\%-4\%=51\%$ fins a $55\%+4\%=59\%$. \textbf{Afirma correctament:} és molt
probable que la satisfacció real de \emph{totes} les famílies estigui entre el $51\%$ i el
$59\%$. \textbf{No} es pot dir que sigui exactament el $55\%$.
\resposta{Interval $[51\%,\,59\%]$; el valor real és molt probablement dins d'aquest rang.}
\end{exresolt}

\begin{exresolt}[detectar una interpretació incorrecta]
Algú llegeix «$48\%\pm 2\%$» i conclou: «exactament el $48\%$ de la població opina X».
Per què és incorrecte? Com ho diries bé?
\solucio
És incorrecte perquè tracta el $48\%$ com un valor \textbf{exacte}, ignorant el marge. La
lectura correcta és: «és molt probable que el valor real estigui \textbf{entre el 46\% i el
50\%}». El resultat d'una mostra sempre s'expressa amb el seu interval, no com una xifra
exacta.
\resposta{Ignora el marge; el correcte és l'interval $[46\%,\,50\%]$.}
\end{exresolt}

\subsection{Exercicis per fer}

\begin{exercicis}
\begin{exlist}
  \item Una enquesta dona «$60\%\pm 3\%$». Construeix l'interval de confiança.
  \item Una enquesta dona «$38\%\pm 5\%$». Entre quins valors és probable que estigui el
        valor real de la població?
  \item Escriu, amb paraules, què afirma correctament una enquesta que diu «el $70\%$ dels
        usuaris recomanaria el servei, amb un marge de $\pm 4\%$».
  \item Digues si aquesta interpretació és correcta o no, i per què: «$45\%\pm 2\%$ vol dir
        que exactament el $45\%$ opina X».
  \item \textbf{(Error típic)} Algú diu que el «$\pm 3\%$» significa que «el $3\%$ de la
        gent ha respost malament». Explica per què és fals i què vol dir realment el marge.
  \item \textbf{(Interpretació)} Dues enquestes donen «$50\%\pm 3\%$» i «$50\%\pm 8\%$».
        Totes dues tenen el mateix resultat central. Quina és més \emph{precisa}? Per què?
\end{exlist}
\end{exercicis}

% ---------------------------------------------------------------------
\fullsolucions{Sessió 5}
\begin{solucions}
\begin{sollist}
  \item $60\%-3\%=57\%$ fins a $60\%+3\%=63\%$: interval $[57\%,\,63\%]$.
  \item $38\%-5\%=33\%$ fins a $38\%+5\%=43\%$: entre el $33\%$ i el $43\%$.
  \item Resposta oberta. Idea: és molt probable que, a tota la població d'usuaris, la
        proporció que recomanaria el servei estigui entre el $66\%$ i el $74\%$; no que
        sigui exactament el $70\%$.
  \item Incorrecta: ignora el marge i tracta el $45\%$ com a exacte. El correcte és que el
        valor real està molt probablement entre el $43\%$ i el $47\%$.
  \item Resposta oberta. Idea: el marge no té res a veure amb respostes errònies; expressa
        que, com que s'ha preguntat a una mostra i no a tothom, el valor real pot diferir
        del mesurat fins a un $3\%$ amunt o avall (incertesa del mostreig).
  \item La de $\pm 3\%$ és més precisa: el seu interval ($[47\%,53\%]$) és més estret que
        el de $\pm 8\%$ ($[42\%,58\%]$). Menys marge vol dir una estimació més precisa.
\end{sollist}
\end{solucions}
